\chapter{KMS-совместимый канал выхода и граница детального баланса}
\label{ch:v10-breathing-k43-kms-output-channel}

\section{Канонический двухуровневый канал}

На подносителе $W_Y=\operatorname{span}\{|0\rangle,|Y\rangle\}$ зададим
вероятности нисходящего и восходящего переходов
\begin{equation}
 p_\downarrow=\frac16,\qquad p_\uparrow=\frac1{12},
 \qquad \frac{p_\uparrow}{p_\downarrow}
 =\frac12=e^{-\log2}.
 \label{eq:v10-k43-kms-rates}
\end{equation}
Операторы Крауса можно взять в виде
\begin{equation}
 K_0=\begin{pmatrix}\sqrt{11/12}&0\\0&\sqrt{5/6}\end{pmatrix},
 \quad
 K_\downarrow=\begin{pmatrix}0&\sqrt{1/6}\\0&0\end{pmatrix},
 \quad
 K_\uparrow=\begin{pmatrix}0&0\\\sqrt{1/12}&0\end{pmatrix}.
 \label{eq:v10-k43-kms-kraus}
\end{equation}
Они удовлетворяют $\sum_iK_i^\dagger K_i=I_2$, поэтому определяют полностью
положительный сохраняющий след канал. Потеря чистого возбуждённого состояния
равна $p_\downarrow=1/6$ и точно совпадает с каноническим спектральным
свидетелем $R_{\rm out}$.

\section{Конечная KMS-температура}

Матрица переходов популяций равна
\begin{equation}
 T_{\rm KMS}=
 \begin{pmatrix}
 11/12&1/6\\
 1/12&5/6
 \end{pmatrix}.
 \label{eq:v10-k43-kms-transition}
\end{equation}
Она фиксирует Gibbs-вектор
\begin{equation}
 \pi_{\rm KMS}=\begin{pmatrix}2/3\\1/3\end{pmatrix},
 \qquad
 \rho_{\rm KMS}=\operatorname{diag}(2/3,1/3),
 \label{eq:v10-k43-kms-state}
\end{equation}
соответствующий $\beta\Delta=\log2$. Спектр $T_{\rm KMS}$ равен
$(1,3/4)$.

Однако равновесные направленные потоки совпадают:
\begin{equation}
 J_\downarrow=\frac13\frac16=\frac1{18},
 \qquad
 J_\uparrow=\frac23\frac1{12}=\frac1{18},
 \qquad
 J_{\rm net}=0.
 \label{eq:v10-k43-kms-zero-current}
\end{equation}
Таким образом, конечный KMS-канал реализует вероятность потери $1/6$, но
не создаёт стационарный сквозной выход.

\section{Нулевая температура}

Если положить $p_\uparrow=0$, остаётся обычный амплитудно-затухающий канал
с той же нисходящей вероятностью $1/6$. Его стационарное состояние равно
$|0\rangle\langle0|$ и не является верным конечным KMS-состоянием.
Следовательно, односторонний выход достигается ценой перехода к граничному
нулетемпературному резервуару.

\section{Условный родитель}

Условия полной положительности $u_{\rm CP}$, KMS-стационарности
$u_{\rm KMS}$ и спектрального совпадения $u_{\rm match}$ объединяются в
\begin{equation}
 \mathcal P_{\rm KMS,out}
 =\frac12(u_{\rm CP}-1)^2
 +\frac12(u_{\rm KMS}-u_{\rm CP})^2
 +\frac12(u_{\rm match}-u_{\rm KMS})^2.
 \label{eq:v10-k43-kms-parent}
\end{equation}
Гессиан имеет ранг $3$, определитель $1$ и ведущие миноры $(2,3,1)$.
Канал безразмерен и не добавляет новой строки в размерную карту: её
ранг/ядро остаются $3/1$.

\begin{theorem}[Канальная реализуемость одной шестой]
Существует точный полностью положительный сохраняющий след канал, у которого
нисходящая вероятность равна $1/6$, KMS-отношение скоростей равно $1/2$ и
стационарное состояние равно $\operatorname{diag}(2/3,1/3)$.
\end{theorem}

\begin{nogocriterion}[Детальный баланс запрещает сквозной ток]
В стационарном конечном KMS-состоянии прямой и обратный потоки равны
$1/18$, поэтому чистый выход равен нулю. Поддерживаемое дыхание требует
неравновесного привода, двух резервуаров с различными сродствами или явного
нарушения детального баланса.
\end{nogocriterion}

\begin{protocol}[Статус KMS-канала]\begin{enumerate}
\item условная архитектура: \textbf{$10/10$};
\item условное замыкание: \textbf{$8/8$};
\item точный CPTP-канал: \textbf{$1/1$};
\item совпадение потери с $1/6$: \textbf{$1/1$};
\item стационарный сквозной ток при KMS: \textbf{$0/1$};
\item происхождение неравновесного привода и ванн: \textbf{$0/2$};
\item абсолютный масштаб: \textbf{$0/1$};
\item следующий гейт: \textbf{двухрезервуарный неравновесный ток}.
\end{enumerate}\end{protocol}

Машинный сертификат сохранён в
\path{s2t_v10_cell_birth_four_volume_induced_newton_breathing_anomaly_k43_kms_output_channel_parent_origin_gate_results.json}.